(1) 设 \(n_A=a_0-\xi,n_B=b_0-\xi,n_C=c_0+\xi\)，均正，\(p_i=n_iRT/V\)。理想混合物化学势
\(\mu_i=\mu_i^\circ(T)+RT\ln(p_i/p^\circ)\)，反应自由能导数
\[
{dF\over d\xi}=\mu_C-\mu_A-\mu_B
=RT\ln(Q_p/K_p),\quad
Q_p={(p_C/p^\circ)\over(p_A/p^\circ)(p_B/p^\circ)}.
\]
平衡满足 \(Q_p=K_p\)。
(2) 进一步
\[
{d^2F\over d\xi^2}=RT\left({1\over n_C}+{1\over n_A}+{1\over n_B}\right)>0,
\]
故在正物质的量区间严格凸。若 \(a_0=b_0=n,c_0=0\)，则
\[
Q_p={\xi(Vp^\circ/RT)\over(n-\xi)^2}=K_p,\quad 0<\xi<n .
\]
左侧随 \(\xi\) 增大由0趋无穷，解唯一。
(3) Gibbs–Helmholtz关系 \(d(\Delta G^\circ/T)/dT=-\Delta H^\circ/T^2\)，结合 \(\Delta G^\circ=-RT\ln K_p\)，得
\[
{d\ln K_p\over dT}={\Delta H^\circ\over RT^2}.
\]
放热 \(\Delta H^\circ<0\) 则升温降低K，平衡偏反应物。
(4) 定温定容自发方向须使 \(dF/d\xi\,d\xi\le0\)。若 \(Q_p<K_p\)，则 \(dF/d\xi<0\) 且正向反应使F下降；反之逆向。体积增大使 \(Q_p\propto V\) 增大，平衡向气体总摩尔数较多的左侧移动；引入活度系数后不可再用理想分压比取代热力学活度。
